\section{Hell Is Eternal Punishment}

Hell is the everlasting punishment of those who die in mortal sin without repenting and accepting God's mercy. Its chief punishment is the loss of God: exclusion forever from the blessed vision and communion for which we were created. The damned also suffer the punishment Christ describes as everlasting fire. The souls of the damned suffer immediately after death; at the resurrection, their bodies too will share their punishment (CCC 1033--1038).

Christ himself declares the sentence: ``Depart from me, you cursed, into everlasting fire'' (Mt 25:41). These words belong to the Gospel of our salvation. The Lord who speaks them became man for us, sought sinners, forgave sins, and gave his life upon the Cross. His warning proceeds from his knowledge of the ruin from which he came to deliver us. Faith receives the warning from the same Lord whose promises it trusts.

The sentence is a judgment. God truly judges and punishes sin. The sinner's free rejection of grace and the Judge's just sentence belong to the same reality. Damnation does not result from a divine decision to create someone for misery: God predestines no one to hell. It requires a willful turning away from him and persistence in that rejection until death (CCC 1037). Neither does it require the sinner to desire punishment for its own sake. A person may know that a choice gravely offends God and still choose it for the satisfaction, revenge, power, or advantage it promises. The chosen evil can be loved while its eternal consequence is pushed from thought.

God created the whole person for communion with himself. The body shares in our response to his grace: we act through it, and he will raise it and judge what we have done. Everlasting blessedness is the fulfillment of that created life. Its final loss therefore reaches beyond the deprivation of any particular pleasure or possession. A good which sustains us during earthly life can never replace the end for which we exist.

Hell also differs utterly from purgatory. Those undergoing purification belong to Christ and will enter heaven. The damned are deprived of salvation; their punishment does not prepare them for the blessed vision. Nor do they cease to exist. At the end of the judgment scene Christ sets everlasting punishment beside life everlasting (Mt 25:46). Both are final destinies. The earthly interval in which a person can turn from sin and accept grace comes to an end at death (CCC 1021, 1030--1032).

We meditate upon this truth while repentance remains possible. The sinner who has spent years refusing a known duty can still turn to Christ; the passage of those years has not made his present refusal necessary. This possibility is itself mercy. Another day has been given, the Gospel remains before us, and the Church continues to minister Christ's forgiveness. Yet the time received calls for a response now. Continuing to postpone conversion uses the very gift that makes conversion possible to prolong the sin from which we need deliverance.

As the year draws toward its close, an examination of conscience can expose how little a defended fault has changed despite everything else that happened. The occasions were different, but the excuse remained available. We may have become quicker at explaining why a duty could wait. We may also discover that grace has weakened an attachment which once seemed immovable. Such a review concerns the state of the soul before God. Outward activity alone cannot tell us whether we have grown in charity or become more practiced in avoiding its demands.

The saints considered hell in order to live and die in God's friendship. Their meditations examine how sin estranges us from him and how his grace restores us. They also compare the good of remaining faithful with the suffering fidelity may require. Christ's warning gives a reason to bear that suffering: the threatened loss is greater than the temporal good we are tempted to preserve by sin.

